% Shared preamble
\usepackage[a4paper,left=24mm,right=24mm,top=24mm,bottom=24mm,headheight=16pt,headsep=8mm]{geometry}
\usepackage{mathtools,amssymb,amsthm}
\usepackage[protrusion=false,expansion=false]{microtype}
\usepackage{booktabs,array}
\usepackage[font=small,labelfont=bf,labelsep=period]{caption}
\usepackage[shortlabels]{enumitem}
\usepackage{needspace}
\usepackage{tikz}
\usetikzlibrary{calc,positioning,arrows.meta}
\definecolor{ink}{HTML}{213650}
\definecolor{sea}{HTML}{19706F}
\definecolor{rust}{HTML}{AA5034}
\usepackage[numbers,sort&compress]{natbib}
\renewcommand\bibfont{\small}
\usepackage[nottoc,notlof,notlot]{tocbibind}
\usepackage{fancyhdr}
\usepackage{url}
\usepackage[colorlinks,linkcolor=ink,citecolor=sea,urlcolor=ink,psdextra]{hyperref}
\usepackage{bookmark}

\fancyhf{}
\newcommand\shorttitle{OpenAI}
\fancyhead[L]{\small\shorttitle}
\fancyhead[R]{\small\thepage}
\renewcommand\headrulewidth{0pt}
\pagestyle{fancy}
\newcommand{\articletitle}[2]{\title{\fontsize{20}{24}\selectfont\bfseries\color{ink}#1}\renewcommand\shorttitle{#2}%
  \hypersetup{pdftitle={#1},pdfauthor={OpenAI}}}
\author{OpenAI\quad 著\qquad 中文译本}
\date{2026 年 11 月 24 日\\[.35em]{\footnotesize\normalfont\href{https://github.com/openai/math/blob/3344da59991ed94f9ba711217b29c1d760badfa6/preprints/The-matrix-multiplication-exponent-is-nine-fourths-November-24-2026/paper.pdf}{英文原文}\quad$\cdot$\quad OpenAI Math 预印本}}
\setcounter{tocdepth}{1}

\newtheorem{theorem}{定理}[section]
\newtheorem{lemma}[theorem]{引理}
\newtheorem{claim}[theorem]{断言}
\newtheorem{corollary}[theorem]{推论}
\newtheorem{proposition}[theorem]{命题}
\theoremstyle{definition}
\newtheorem{definition}[theorem]{定义}
\newtheorem{remark}[theorem]{注}
\newtheorem{example}[theorem]{例}
\numberwithin{equation}{section}
\emergencystretch=3em
\clubpenalty=3500
\widowpenalty=3500
\raggedbottom

% For numbers of results in companion papers
\makeatletter
\newcommand\ext[2]{\@ifundefined{ref@#1:#2}{??\@latex@warning{Undefined companion label #1:#2}}{\csname ref@#1:#2\endcsname}}
\newcommand\extdef[3]{\@namedef{ref@#1:#2}{#3}}
\makeatother
\extdef{G}{app:trace}{16.2}
\extdef{G}{def:ports}{11.1}
\extdef{G}{def:tag}{8.1}
\extdef{G}{lem:slack-compaction}{10.2}
\extdef{G}{lem:spectral-comparison}{2.2}
\extdef{G}{prop:joint-convexity}{3.3}
\extdef{G}{prop:oracle-descent}{14.1}
\extdef{G}{sec:limits}{2.3}
\extdef{G}{sec:protocol}{9}
\extdef{G}{sec:spectra}{2}
\extdef{G}{thm:crossing}{4.1}
\extdef{G}{thm:extension}{3.1}
\extdef{G}{thm:geometry}{4.3}
\extdef{P}{cor:graph-entropy}{5.2}
\extdef{P}{prod:main}{6.1}
\extdef{P}{thm:batches}{12.1}



\newcommand\titlebreak{\texorpdfstring{\\}{}}
\settocbibname{参考文献}

% Lightweight contents, with numbered bookmarks
\bookmarksetup{numbered}
\makeatletter
\renewcommand\@pnumwidth{2.1em}
\renewcommand\@dotsep{12}
\renewcommand*\l@section[2]{\addvspace{.22em}\@dottedtocline{1}{\z@}{2em}{#1\nobreak\hspace{.5em}}{#2}}
\makeatother


\renewcommand{\proofname}{证明}
\renewcommand{\abstractname}{摘要}
\renewcommand{\contentsname}{目录}
\renewcommand{\figurename}{图}
\renewcommand{\tablename}{表}
\renewcommand{\partname}{部分}
\renewcommand{\refname}{参考文献}
\ctexset{part/name={第,部分},part/number=\chinese{part}}
\setlength{\parindent}{2em}
